\chapter{Conexión y Recorridos}

\section{Trayectorias y Caminos}
Un camino en un grafo es una secuencia de vértices donde cada par consecutivo está conectado por una arista.

\subsection{Tipos de caminos}
\begin{itemize}
  \item \textbf{Camino simple}: No repite vértices.
  \item \textbf{Ciclo}: Camino que comienza y termina en el mismo vértice.
  \item \textbf{Trayectoria}: Camino sin aristas repetidas.
\end{itemize}

\section{Grafos Conexos}
Un grafo es conexo si existe un camino entre cualquier par de vértices.

\subsection{Componentes conexas}
Una componente conexa es un subgrafo maximal que es conexo.

\section{Recorridos en Grafos}

\subsection{Búsqueda en profundidad (DFS)}

\begin{algorithm}[H]
\caption{DFS (Búsqueda en Profundidad)}
\label{alg:dfs}
\begin{algorithmic}[1]
\Require Un grafo $G=(V,E)$ y un vértice inicial $s$
\Ensure Recorrido DFS
\State visitado$[v] \gets \text{falso}$ para todo $v \in V$
\Procedure{DFS}{$u$}
  \State visitado$[u] \gets \text{verdadero}$
  \For{cada $v \in \text{adyacentes}(u)$}
    \If{$\neg$visitado$[v]$}
      \State DFS$(v)$
    \EndIf
  \EndFor
\EndProcedure
\State DFS$(s)$
\end{algorithmic}
\end{algorithm}

\subsection{Búsqueda en anchura (BFS)}

\begin{algorithm}[H]
\caption{BFS (Búsqueda en Anchura)}
\label{alg:bfs}
\begin{algorithmic}[1]
\Require Un grafo $G=(V,E)$ y un vértice inicial $s$
\State visitado$[v] \gets \text{falso}$ para todo $v \in V$
\State visitado$[s] \gets \text{verdadero}$
\State cola $\gets [s]$
\While{cola $\neq \emptyset$}
  \State $u \gets \text{POPFRONT}(cola)$
  \For{cada $v \in \text{adyacentes}(u)$}
    \If{$\neg$visitado$[v]$}
      \State visitado$[v] \gets \text{verdadero}$
      \State cola $\gets$ cola $+$ $[v]$
    \EndIf
  \EndFor
\EndWhile
\end{algorithmic}
\end{algorithm}

\section{Caminos Eulerianos y Hamiltonianos}

\subsection{Caminos Eulerianos}
Un camino euleriano recorre todas las aristas exactamente una vez. Un grafo conexo tiene un ciclo euleriano si y solo si todos los vértices tienen grado par.

\subsection{Caminos Hamiltonianos}
Un camino hamiltoniano visita todos los vértices exactamente una vez. Teorema de Dirac: Si un grafo con $n \geq 3$ vértices tiene $\deg(v) \geq n/2$ para todo $v$, entonces tiene un ciclo hamiltoniano.

\section{Implementación en Haskell}

\begin{lstlisting}[caption={Algoritmos de recorrido en Haskell}, label=lst:grafos-haskell]
-- DFS
dfs :: Eq a => a -> [(a, [a])] -> [a]
dfs inicio grafo = dfsAux inicio []
  where
    dfsAux v visitados
      | v `elem` visitados = visitados
      | otherwise = foldl (\acc w -> dfsAux w acc) (v:visitados) (adyacentes v grafo)

-- BFS
bfs :: Eq a => a -> [(a, [a])] -> [a]
bfs inicio grafo = bfsAux [inicio] []
  where
    bfsAux [] visitados = visitados
    bfsAux (v:cola) visitados
      | v `elem` visitados = bfsAux cola visitados
      | otherwise = bfsAux (cola ++ adyacentes v grafo) (visitados ++ [v])

-- Verificar si es conexo
esConexo :: Eq a => [(a, [a])] -> Bool
esConexo grafo = all (`elem` reachable) (map fst grafo)
  where (inicio:_) = map fst grafo
        reachable = dfs inicio grafo

-- Verificar si tiene ciclo euleriano
cicloEuleriano :: Eq a => [(a, [a])] -> Bool
cicloEuleriano grafo = esConexo grafo && all even (map (length . snd) grafo)
\end{lstlisting}

\section{Ejercicios}
\begin{enumerate}
  \item Ejecutar DFS y BFS en un grafo de ejemplo.
  \item Demostrar que un grafo tiene un ciclo euleriano si todos los vértices tienen grado par.
  \item Implementar en Haskell un algoritmo que encuentre un camino euleriano.
\end{enumerate}